\begin{musica}{Muñequita Linda}{María Grever}{Tuna de Engenharia da Universidade do Porto}
Te quiero \ac{C Am}{}_ __ dijiste \ac{C Am}{}_ __
To\ac{F}{man}do mis \ac{Am7}{ma}nos
En\ac{F}{tre} tus mani\ac{F#dim7}{tas} de \ac{Gm}{bian}co mar\ac{Gm7 C}{fil} _
\ac{A7}{Y} senti en mi \ac{Dm7}{pe}cho \ac{F}{}_ un \ac{G7}{fuer}te la\ac{C}{ti}do
De\ac{A7}{pués} un sus\ac{Dm7}{pi}ro
\ac{F}{Y} luego el chas\ac{G}{qui}to de un \ac{G13b}{be}so fe\ac{C}{bril}. \ac{C7}{}_

Muñequita \ac{F}{lin}da \ac{Dm}{}__ de cabellos de o\ac{F}{ro}, \ac{Dm}{}__
De dientes de per\ac{Gm C7}{la}, __ labios de ru\ac{F}{bi}. \ac{C7}{}__
Dime si me \ac{F}{quie}res \ac{Dm}{}__ como yo te a\ac{F}{do}ro \ac{Dm}{}__
Si de mi te a\ac{Gm}{cuer}das \ac{C7}{}__ como yo de \ac{F}{ti}. \ac{F7}{}__

As veces es\ac{A#}{cu}cho \ac{Bdim}{}____ un \ac{E7}{eco} di\ac{Am}{vi}no \ac{D7}{}__
Que envuelto en la \ac{Gm}{bri}sa \ac{G7}{}__ parece de\ac{C}{cir}. \ac{C7}{}__

Si te quiero \ac{F}{mu}cho, \ac{Dm}{}__ mucho, mucho, \ac{F}{mu}cho, \ac{Dm}{}__
Tanto como enton\ac{Gm}{ces} siempre has\ac{C7}{ta} mo\ac{F}{rir}. \ac{A#m F}{}___ _

\posicoes{C7: x32010, Dm: xx0231, Dm7: xx0211, F: 133211, F7: 131211,
  F#dim7: xx4545, G13b: 3x044x, G7: 320001, Gm: 355333, Gm7: 353333,
  A7: x02020, Am: x02210, Am7: x02010, A#: x13331, A#m: x13321,
  Bdim: x23131}
\end{musica}
